\documentclass[11pt]{article}
\usepackage[margin=1.1in]{geometry}
\usepackage{amsmath,amssymb,amsthm}
\usepackage[colorlinks=true,linkcolor=blue,citecolor=blue]{hyperref}

\newtheorem{theorem}{Theorem}[section]
\newtheorem{lemma}[theorem]{Lemma}
\newtheorem{corollary}[theorem]{Corollary}
\newtheorem{proposition}[theorem]{Proposition}
\theoremstyle{definition}
\newtheorem{definition}[theorem]{Definition}
\theoremstyle{remark}
\newtheorem{remark}[theorem]{Remark}
\newtheorem{question}[theorem]{Question}

\newtheorem*{thmA}{Theorem A}
\newtheorem*{thmB}{Theorem B}
\newtheorem*{thmC}{Theorem C}
\newtheorem*{thmD}{Theorem D}

\newcommand{\Z}{\mathbb{Z}}
\newcommand{\R}{\mathbb{R}}
\newcommand{\ord}{\operatorname{ord}}
\newcommand{\bitlen}{\operatorname{len}}
\newcommand{\LD}{\operatorname{LD}}
\DeclareMathOperator{\vv}{v}

\title{No-Go Theorems for One-Step Lyapunov Potentials\\ for the $3x+1$ Map}
\author{[author]\thanks{Repository: \texttt{github.com/docbgm2002/collatz-things}.
Every numbered claim is machine-verified in exact integer or rational
arithmetic; see Section~\ref{sec:verification}.}}
\date{Draft of July 2, 2026}

\begin{document}
\maketitle

\begin{abstract}
A natural strategy for the $3x+1$ (Collatz) problem is to exhibit a
potential $\Phi(x)=\log_2x+g(x)$ that is nonincreasing along the
accelerated odd map $f(x)=(3x+1)/2^{\vv_2(3x+1)}$, with a correction $g$
computed from simple features of $x$. It is folklore that no such
monotone invariant is known; we prove that for broad natural classes of
corrections none can exist. The master mechanism is \emph{shadowing an
expanding rational cycle}: the cycle $-5\mapsto-7\mapsto-5$ has
multiplier $3^2/2^3>1$, and its positive-integer shadows
$x\equiv-5\pmod{2^N}$ traverse its coordinate itinerary while gaining
value, so any correction that cannot separate $x$ from
$f^2(x)=\tfrac98x+\tfrac58$ is refuted by two integers --- the smallest
certificate being $1275\mapsto1913\mapsto1435$. This excludes all
corrections $g\bigl(x\bmod2^m,\ \tau(x),\ \bitlen(x),\
\lambda_1(x),\dots,\lambda_d(x)\bigr)$, where $\tau$ is the
trailing-ones count and the $\lambda_i$ are suffix-pattern detectors of
arbitrary depth (Theorem~A). Iterating the shadow and applying
Dirichlet approximation to $\log_2\frac98$ extends the exclusion to
corrections using any fixed number of leading binary digits
(Theorem~B); a Mersenne burn argument excludes all bounded corrections
of arbitrary dependence (Theorem~C). The detectors are not ad hoc: they
read a canonical structure, an exact ancestry tower
$w_d(M)=\bigl(2^{M+2d}-2^{2d+1}+3^{d}\bigr)/3^{d}$ with
$f^{d}(w_d(M))=2^{M}-1$, supported on $M\equiv1\pmod{2\cdot3^{d-1}}$
with a carry-free periodic binary normal form (Theorem~D). The
surviving one-step class reads the quantized full logarithm, where
further no-go progress becomes conjecture-adjacent.
\end{abstract}

\section{Introduction}

For odd $x\in\Z_{>0}$ let
\[
f(x)=\frac{3x+1}{2^{\vv_2(3x+1)}},
\]
the accelerated (odd-to-odd) form of the Collatz map; the $3x+1$
conjecture asserts that every orbit reaches $1$. Among the strategies
proposed for the problem, one of the most natural is the
\emph{potential-function} (Lyapunov) method: find
$\Phi:2\Z+1\to\R$ with $\Phi(f(x))\le\Phi(x)$ for all odd $x>1$,
ideally of the coercive form $\Phi(x)=\log_2x+g(x)$ with $g$ a
correction of lower order, so that monotonicity constrains orbit
growth. Heuristics support the idea: a ``random'' odd step multiplies
$x$ by $3/2^{e}$ with $e$ geometrically distributed of mean $2$, so
$\log_2x$ decreases by $\log_2\frac43$ \emph{on
average}~\cite{Lagarias1985}. The difficulty, as surveys note, is that
no correction making the decrease pointwise has ever been
found~\cite{Lagarias2010}. To the best of our knowledge the literature
contains no theorem \emph{excluding} a concrete class of such
potentials; the absence is reported as an empirical fact. This paper
proves exclusions, and locates the precise boundary at which the
exclusion program must stop.

Throughout, $\tau(x)=\vv_2(x+1)$ is the number of trailing
\texttt{1}s in the binary expansion of odd $x$ (the \emph{fuel}:
Lemma~\ref{lem:fuse} shows an odd step converts one unit of fuel into a
multiplication by almost exactly $\tfrac32$); $\bitlen(x)$ is the bit
length; $\LD_j(x)=\lfloor x/2^{\bitlen(x)-j}\rfloor$ is the top $j$
bits; and $\lambda_1,\lambda_2,\dots$ are the suffix-pattern detectors
of Definition~\ref{def:detector}, which read the run length of a
canonical hierarchy of periodic binary patterns (levels $1$ and $2$
being $(\texttt{10})^{k}\texttt{1001}$ and
$(\texttt{111000})^{k}\texttt{01}$). Our main impossibility results:

\begin{thmA}[master no-go; Theorem~\ref{thm:A}]
For every $m,d\ge0$ and every function $g$, the potential
$\log_2x+g\bigl(x\bmod2^m,\ \tau(x),\ \bitlen(x),\
\lambda_1(x),\dots,\lambda_d(x)\bigr)$ is not nonincreasing along $f$
on odd $x>1$.
\end{thmA}

\begin{thmB}[leading digits; Theorem~\ref{thm:B}]
For every $m,j,d\ge0$ and every $g$, the potential
$\log_2x+g\bigl(x\bmod2^m,\ \tau(x),\ \LD_j(x),\
\lambda_1(x),\dots,\lambda_d(x)\bigr)$ is not nonincreasing along $f$
on odd $x>1$.
\end{thmB}

\begin{thmC}[bounded corrections; Theorem~\ref{thm:C}]
For every bounded $G:2\Z+1\to\R$, of arbitrary dependence on $x$, the
potential $\log_2x+G(x)$ is not nonincreasing along $f$ on odd $x>1$.
\end{thmC}

Theorem~A is proved by two integers per parameter choice: the map has
the expanding rational cycle $-5\mapsto-7\mapsto-5$, of multiplier
$3^2/2^3=\tfrac98>1$ (expansion is exactly why it lives at negative
integers), and its shadow at positive $x\equiv-5\pmod{2^N}$ reproduces
the cycle's coordinate itinerary while the value gains a factor
$>\tfrac98$ per period. Any correction blind to the difference between
$x$ and $f^2(x)=\tfrac98x+\tfrac58$ is refuted outright; the smallest
certificate is $1275\mapsto1913\mapsto1435$. Theorem~B iterates the
shadow $a$ times and chooses $a$, by Dirichlet approximation to the
irrational $\log_2\tfrac98$, so that the accumulated gain sits next to
an exact power of $2$ and the leading digits return. Theorem~C is a
one-paragraph consequence of the Mersenne burn.

The detectors quantified over in Theorems A and B are given meaning by
a structural result of independent interest:

\begin{thmD}[Mersenne ancestry tower; Theorem~\ref{thm:D}]
For $d\ge1$ and $M\equiv1\pmod{2\cdot3^{d-1}}$, $M>1$, the integer
$w_d(M)=\bigl(2^{M+2d}-2^{2d+1}+3^{d}\bigr)/3^{d}$ is odd with
$\tau=1$, satisfies $3w_d(M)+1=4\,w_{d-1}(M)$ and hence
$f^{d}(w_d(M))=2^{M}-1$ with every step of value ratio in
$(\tfrac34,1)$; its binary expansion is exactly periodic (carry-free)
of period $2\cdot3^{d-1}$; and $w_d(M)\equiv1-2^{2d+1}3^{-d}
\pmod{2^{M-1}}$, distinct levels separating $2$-adically at bit
$2\min(d,d')+1$ exactly.
\end{thmD}

The tower is the complete uniform-$e{=}2$ ancestry of the Mersenne
numbers --- the slowest exact approach to maximal fuel --- and
$\lambda_d$ detects membership in its level-$d$ suffix pattern. The
congruence domain is the order of $2$ modulo $3^d$, and the carry-free
normal form is an instance of the lifting-the-exponent lemma.

\subsection*{Method and provenance}

The constraint ``$\Phi$ nonincreasing'' restricted to any coordinate
window is a difference-constraint system on the values of $g$; it is
infeasible if and only if the coordinate graph of real orbit steps
contains a directed cycle whose product of value ratios exceeds $1$
(Section~\ref{sec:certificates}). Mining such cycles by Bellman--Ford
over odd $x\le10^6$ produces startlingly short certificates --- e.g.\
the pair $57\mapsto43$, $27\mapsto41$ already refutes all
$g(x\bmod16,\tau)$ --- and inspection of the mined witnesses revealed
that they shadow the $-5$ cycle, which led to the uniform theorems
above. We record this provenance because the historical route of this
project ran in the opposite direction: Theorems A--C were first proved
for subclasses by designed families (a Mersenne ``burn'' that forces
linear growth of $g$ along a frozen coordinate line, against a
``recharge'' family that caps it), a method that yields an exact price
list --- each level of suffix detection buys a would-be potential
exactly $\log_2\frac43$ of budget against an unbounded liability ---
but is dominated, as a source of impossibility statements, by the
two-integer shadow certificates. The family analysis survives in the
repository accompanying this paper and in the structural
Theorem~D.

\subsection*{The boundary}

A coordinate defeats the shadow only by separating $x$ from
$f^{2a}(x)$ for every good Dirichlet approximation $a$: since the two
agree on all low-order bits, fuel, suffix patterns, and (for suitable
$a$) leading digits, the surviving one-step class reads the
\emph{quantized full logarithm} $\lfloor2^{j}\log_2x\rfloor$. There a
closed certificate needs total gain in $(0,2^{-j})$, which pure shadows
cannot supply; and as $j\to\infty$ the class approaches arbitrary
corrections, for which the existence of a monotone potential is
conjecture-adjacent. Section~\ref{sec:boundary} makes this precise.

\subsection*{Relation to the literature}

The parity-vector structure underlying the heuristic goes back to
Terras~\cite{Terras1976} and Everett~\cite{Everett1977}; density
results culminate in~\cite{KrasikovLagarias2003} and, by different
methods, in Tao's almost-all theorem~\cite{Tao2022}. Cycle exclusions
rest on transcendence bounds~\cite{Steiner1977,SimonsDeWeger2005} and
computation~\cite{Barina2021}. Backward (predecessor) structure of the
Collatz graph is developed at the level of density and measure by
Wirsching~\cite{Wirsching1998}; the tower of Theorem~D is a
distinguished explicit chain in that graph, and its closed form,
congruence domain, and $2$-adic freezing appear to be new. Forward
identities of burn type, attached to the rational cycles at $-1$, $-5$,
$-17$, appear in Andaloro~\cite{Andaloro2000} --- the shadows of
Section~\ref{sec:shadow} are those identities viewed through the
certificate lens. A search of the literature, including the annotated
bibliographies~\cite{LagariasBibI,LagariasBibII}, located no prior
impossibility theorems for concrete classes of one-step potentials; we
would welcome references.

\section{Preliminaries}\label{sec:prelim}

Odd integers are written $x$; $\vv_p$ is the $p$-adic valuation;
$\tau(x)=\vv_2(x+1)$; congruential inverses are taken in
$(\Z/2^m\Z)^{\times}$. Two standard facts are used.

\begin{lemma}[LTE, special case]\label{lem:lte}
For even $n$, $\vv_3(2^{n}-1)=1+\vv_3(n)$; in particular
$\vv_3\bigl(2^{2\cdot3^{d-1}}-1\bigr)=d$.
\end{lemma}

\begin{lemma}\label{lem:order}
$2$ is a primitive root modulo $3^d$ for every $d\ge1$; hence
$\ord_{3^d}(2)=\varphi(3^d)=2\cdot3^{d-1}$.
\end{lemma}

\begin{proof}
$2$ generates $(\Z/9\Z)^{\times}$ (its order divides $6$ and is not
$1,2,3$), and a primitive root mod $9$ is a primitive root mod $3^d$
for all $d$.
\end{proof}

\begin{lemma}[fuse]\label{lem:fuse}
If $\tau(x)=t\ge2$ then $\vv_2(3x+1)=1$, $f(x)=\tfrac{3x+1}{2}$,
$\tau(f(x))=t-1$, and $f(x)/x=\tfrac32+\tfrac1{2x}$.
\end{lemma}

\begin{proof}
Write $x+1=2^{t}u$, $u$ odd. Then
$3x+1=2\,(3\cdot2^{t-1}u-1)$ with odd bracket since $t\ge2$, and
$f(x)+1=3\cdot2^{t-1}u$.
\end{proof}

\begin{remark}\label{rem:andaloro}
The iterated form of Lemma~\ref{lem:fuse} --- that $2^{t}u-1$ ($u$ odd)
iterates in $t$ odd steps to $3^{t}u-1$, the \emph{burn} --- is
classical; it is recorded by Andaloro~\cite{Andaloro2000} (cf.\ the
annotation in~\cite{LagariasBibII}).
\end{remark}

Finally we define the detectors. For $d\ge1$ write
$\ord_d=2\cdot3^{d-1}$ and, for $M\equiv1\pmod{\ord_d}$, $M>1$,
\begin{equation}\label{eq:wd}
w_d(M)=\frac{2^{M+2d}-2^{2d+1}+3^{d}}{3^{d}}
=1+2^{2d+1}\,\frac{2^{M-1}-1}{3^{d}} ,
\end{equation}
an odd integer by Lemma~\ref{lem:order} (Theorem~\ref{thm:D} develops
its structure; nothing beyond the formula and Lemma~\ref{lem:mod32}
below is needed before Section~\ref{sec:tower}).

\begin{definition}\label{def:detector}
For $d\ge1$ let $\beta_d(s)$ be the bit length of
$T^{(d)}_s:=w_d(1+\ord_d\,s)$ and define the \emph{level-$d$ detector}
\[
\lambda_d(x)=\max\bigl\{\,s\ge1:\
x\equiv T^{(d)}_s\ \bigl(\mathrm{mod}\ 2^{\,\beta_d(s)+1}\bigr)\bigr\}
\ \in\ \Z_{\ge1}\cup\{\bot\},
\]
the maximum of the empty set being $\bot$.
\end{definition}

\begin{lemma}\label{lem:mod32}
Every template satisfies $T^{(d)}_s\equiv9\pmod{32}$ if $d=1$ and
$T^{(d)}_s\equiv1\pmod{32}$ if $d\ge2$. Consequently
$\lambda_d(y)=\bot$ for every $d$ whenever
$y\not\equiv1,9\pmod{32}$.
\end{lemma}

\begin{proof}
For $d\ge2$, $2^{2d+1}\equiv0\pmod{32}$ in \eqref{eq:wd}. For $d=1$:
$M$ odd, $M\ge3$, so $2^{M-1}\equiv0\pmod4$, whence
$(2^{M-1}-1)/3\equiv(-1)\cdot3^{-1}\equiv1\pmod4$ and
$w_1(M)\equiv1+8\equiv9\pmod{32}$. Matching any template forces the
mod-$32$ congruence, since every matching modulus
$2^{\beta_d(s)+1}$ contains $32$ ($\beta_d(s)\ge4$).
\end{proof}

\section{The master theorem: shadowing the $-5$ cycle}\label{sec:shadow}

The map has the rational cycle $-5\mapsto-7\mapsto-5$:
$f(-5)=-7$, $f(-7)=-5$, with multiplier $3^2/2^3=\tfrac98>1$. Expanding
cycles live at negative integers (the cycle equation gives
$x=c/(2^{E}-3^{K})$ with $2^{E}<3^{K}$), and their positive shadows
gain value while walking the cycle's coordinate itinerary.

\begin{lemma}[shadow period]\label{lem:S}
Let $N\ge8$, $w$ odd, and $x=2^{N}w-5$. Then
\[
\vv_2(3x+1)=1,\quad f(x)=3\cdot2^{N-1}w-7,\qquad
\vv_2(3f(x)+1)=2,\quad f^{2}(x)=9\cdot2^{N-3}w-5,
\]
so that $8\,f^{2}(x)=9x+5$ and $f^{2}(x)/x=\tfrac98+\tfrac5{8x}$.
Moreover $\tau(x)=\tau(f^2(x))=2$ and $\tau(f(x))=1$;
$x\equiv f^{2}(x)\equiv-5$ and $f(x)\equiv-7\pmod{2^{\,N-3}}$; and
$\lambda_d(x)=\lambda_d(f(x))=\lambda_d(f^2(x))=\bot$ for every $d$.
\end{lemma}

\begin{proof}
Direct computation: $3x+1=2(3\cdot2^{N-1}w-7)$ with odd bracket;
$3f(x)+1=4(9\cdot2^{N-3}w-5)$ with odd bracket; $8f^2(x)=9\cdot2^Nw-40
=9x+5$. Fuel: $x+1=4(2^{N-2}w-1)$;
$f(x)+1=2\cdot3(2^{N-2}w-1)$ with odd second factor;
$f^2(x)+1=4(9\cdot2^{N-5}w-1)$ (odd bracket, $N\ge6$). Residues are
immediate. Detectors: $x\equiv f^2(x)\equiv27$ and
$f(x)\equiv25\pmod{32}$ (using $N-3\ge5$), and
$27,25\notin\{1,9\}$; apply Lemma~\ref{lem:mod32}.
\end{proof}

\begin{theorem}[master no-go]\label{thm:A}
Let $m,d\ge0$ and let $g$ be any real-valued function of
$\bigl(x\bmod2^m,\ \tau(x),\ \bitlen(x),\
\lambda_1(x),\dots,\lambda_d(x)\bigr)$. Then
$\Phi(x)=\log_2x+g(\cdots)$ is not nonincreasing along $f$ on odd
$x>1$.
\end{theorem}

\begin{proof}
Take $N\ge\max(m+3,8)$, $k\ge3$, and $w=2^{k-1}+1$ in
Lemma~\ref{lem:S}. Then additionally
$\bitlen(x)=\bitlen(f^{2}(x))=N+k$, since $w$ places $x$ in the lower
part of its dyadic window: $9(2^{k}+2)<2^{k+4}$. Hence
$x$ and $f^{2}(x)$ have identical coordinate vectors while
$f^{2}(x)>x$, and summing the two monotonicity constraints at $x$ and
$f(x)$ gives
$0\le-\log_2\bigl(f^{2}(x)/x\bigr)<-\log_2\tfrac98<0$.
\end{proof}

The smallest member ($N=8$, $k=3$) is
$1275\mapsto1913\mapsto1435$, with $8\cdot1435=9\cdot1275+5$ and ratio
product $\tfrac{287}{255}$: two integers refute every $g$ on the full
coordinate algebra at $m\le5$.

\begin{remark}[general principle]\label{rem:factories}
Nothing is special about $-5$: every expanding rational cycle of the
map --- every solution of the cycle equation with $2^{E}<3^{K}$, e.g.\
the $K=7$ cycle at $-17$ with multiplier $3^{7}/2^{11}$ --- shadows at
positive integers $\equiv$ (cycle member) $\bmod\ 2^{N}$ into a closed
coordinate itinerary of value gain $3^{K}/2^{E}>1$, refuting every
correction that is $2$-adically local along it. Expanding negative
cycles are certificate factories against local potentials.
\end{remark}

\section{Iterating the shadow: leading digits}\label{sec:leading}

A coordinate can defeat a single shadow period only by separating $x$
from $f^{2}(x)=\tfrac98x+\tfrac58$; the leading digits do (the ratio
$\tfrac98$ shifts them). Iteration removes this refuge.

\begin{theorem}[leading digits do not survive]\label{thm:B}
For every $m,j,d\ge0$ and every $g$, the potential
$\log_2x+g\bigl(x\bmod2^m,\ \tau(x),\ \LD_j(x),\
\lambda_1(x),\dots,\lambda_d(x)\bigr)$ is not nonincreasing along $f$
on odd $x>1$.
\end{theorem}

\begin{proof}
For $x\equiv-5\pmod{2^{\,3a+c}}$, $c\ge5$, Lemma~\ref{lem:S} applies
$a$ times (each period consumes three bits of shadow depth), the
itinerary alternating residues $-5,-7$ with fuel $2,1$ and every
detector $\bot$ throughout, giving the exact closed form
\[
f^{2a}(x)=\frac{9^{a}x+5\,(9^{a}-8^{a})}{8^{a}}
=\Bigl(\tfrac98\Bigr)^{a}x\,(1+\delta),\qquad 0<\delta<\tfrac5x .
\]
Now $\beta=\log_2\tfrac98$ is irrational ($2^{\,p+3q}=3^{\,2q}$ is
impossible: $\vv_3$ of the sides differ), so by Dirichlet there is
$a\ge1$ with $\|a\beta\|<2^{-j-3}$ (distance to the nearest integer
$P$; note $a\beta>0$ regardless of the side). Choose
$N=3a+\max(m,5)+j+8$ and a free odd multiplier $w$ placing the
fractional part of $\log_2x$, $x=2^{N}w-5$, at distance $>2^{-j-2}$
from every breakpoint of $\LD_j$ (possible since
$\log_2(2^{N}w-5)\approx N+\log_2w$ and $\{\log_2w\}$ is dense), with
$N$ large enough that $\log_2(1+\delta)<2^{-j-3}$. Then
$\log_2f^{2a}(x)=\log_2x+P\pm\|a\beta\|+\log_2(1+\delta)$, so
$\LD_j$ closes; all other coordinates close by the itinerary; and
telescoping monotonicity over the $2a$-step walk gives
$0\ge a\beta+\log_2(1+\delta)>0$.
\end{proof}

A concrete instance ($j=4$): $a=53$ gives $53\beta=9.0060\ldots$; the
verifier realizes the full $106$-step walk in exact integers at depth
$N=177$, with top-$4$ bits $\texttt{1011}\to\texttt{1011}$, all
coordinates frozen, and value gain $2^{9.006}$.

\section{Bounded corrections}\label{sec:bounded}

\begin{theorem}[no bounded correction]\label{thm:C}
Let $G:2\Z+1\to\R$ be any bounded function. Then $\log_2x+G(x)$ is not
nonincreasing along $f$ on odd $x>1$.
\end{theorem}

\begin{proof}
Let $\Omega=\sup G-\inf G$. Telescoping monotonicity over the burn of
$2^{M}-1$ (Lemma~\ref{lem:fuse} iterated: $M-1$ steps, each of ratio
$>\tfrac32$) forces $(M-1)\log_2\tfrac32\le\Omega$, false for
$M>1+\Omega/\log_2\tfrac32$.
\end{proof}

Theorems A--C partition the correction landscape: bounded corrections
fail on the burn alone; unbounded corrections through residues, fuel,
length, leading digits, and suffix detectors of every depth fail on
the shadow. Section~\ref{sec:boundary} describes what is left.

\section{The Mersenne ancestry tower}\label{sec:tower}

We now develop the structure behind the detectors. Define
$w_0(M)=2^{M}-1$ and $w_d(M)$ as in \eqref{eq:wd}.

\begin{theorem}\label{thm:D}
For $d\ge1$ and $M\equiv1\pmod{\ord_d}$, $M>1$:
\begin{enumerate}
\item[(a)] \textup{(Domain.)} $3^{d}$ divides
$2^{M+2d}-2^{2d+1}+3^{d}$ iff $M\equiv1\pmod{\ord_d}$.
\item[(b)] \textup{(Form.)} $w_d(M)$ is an odd integer $>1$ with
$\tau(w_d(M))=1$.
\item[(c)] \textup{(Descent.)} $3\,w_d(M)+1=4\,w_{d-1}(M)$; hence
$\vv_2(3w_d(M)+1)=2$, $f(w_d(M))=w_{d-1}(M)$, and
$f^{d}(w_d(M))=2^{M}-1$, each step of value ratio
$\tfrac34+\tfrac1{4w_d}\in(\tfrac34,1)$.
\item[(d)] \textup{(Normal form.)} $B_d:=(2^{\ord_d}-1)/3^{d}$ is an
integer in $(0,2^{\ord_d})$ by Lemma~\ref{lem:lte}, and for
$M=1+\ord_d\,s$,
\[
w_d(M)=1+2^{2d+1}\cdot\bigl[\text{the $\ord_d$-bit block of
}B_d\bigr]^{s},
\]
the bracket denoting $s$ carry-free repetitions: an exactly periodic
binary word of period $2\cdot3^{d-1}$.
\item[(e)] \textup{(Freeze and separation.)}
$w_d(M)\equiv\Lambda_d:=1-2^{2d+1}3^{-d}\pmod{2^{\,M-1}}$; and for
$d\ne d'$, $\vv_2(\Lambda_d-\Lambda_{d'})=2\min(d,d')+1$.
\end{enumerate}
\end{theorem}

\begin{proof}
(a) Modulo $3^{d}$ the numerator is $2^{2d}(2^{M}-2)$, so divisibility
is $2^{M-1}\equiv1\pmod{3^d}$; apply Lemma~\ref{lem:order}.
(b) The second form in \eqref{eq:wd} is $1$ plus an even number, and
$w_d+1=2(1+2^{2d}q)$ with $2^{2d}q$ even.
(c) Clearing $3^{d-1}$, the identity reduces to
$3^{d}+3^{d-1}=4\cdot3^{d-1}$; $w_{d-1}$ is odd, so $\vv_2=2$ exactly;
iterate.
(d) With $M-1=\ord_d\,s$,
$(2^{M-1}-1)/3^{d}=B_d\sum_{i=0}^{s-1}2^{\ord_d i}$, a
base-$2^{\ord_d}$ number all of whose digits equal $B_d<2^{\ord_d}$: no
carries; the $+1$ and the shift by $2^{2d+1}$ cannot interact.
(e) $3^{d}\cdot\tfrac{2^{M-1}-1}{3^{d}}\equiv-1\pmod{2^{M-1}}$; and for
$d>d'$,
$\Lambda_d-\Lambda_{d'}=2^{2d'+1}3^{-d}\bigl(3^{\,d-d'}-2^{\,2(d-d')}\bigr)$
with odd bracket.
\end{proof}

Levels $0,1,2$ are the Mersennes, the family
$x'_M=(2^{M+2}-5)/3=(\texttt{10})^{(M-3)/2}\texttt{1001}$ with
$f(x'_M)=2^M-1$, and $(\texttt{111000})^{k}\texttt{01}$; the tower is
the complete uniform-$e{=}2$ ancestry of maximal fuel, and every tower
member lies on the residue class $8y+1$. Composing the tower with the
burn (Remark~\ref{rem:andaloro}) yields, for each $(d,M)$, an exactly
known orbit segment of length $d+M+1$ terminating on the base-$3$
repunit $(3^{M}-1)/2$.

\begin{lemma}\label{lem:detfacts}
(1) $\lambda_d(T^{(d)}_s)=s$: the detector reads its own tower index,
unboundedly. (2) For $d\ne d'$, with $\nu=2\min(d,d')+1$: for all
$M>\nu+1$ in the level-$d$ progression,
$\lambda_{d'}(w_d(M))\le\nu/\ord_{d'}$.
\end{lemma}

\begin{proof}
(1) At index $s$ the congruence is an equality that holds; at $s'>s$
the modulus exceeds both operands and templates strictly increase.
(2) Suppose $w_d(M)$ matches $T^{(d')}_s$ with $M'=1+\ord_{d'}s$ and,
toward a contradiction, $M'-1\ge\nu+1$. The matching modulus refines
$2^{\nu+1}$ (as $\beta_{d'}(s)\ge M'$), so Theorem~\ref{thm:D}(e) on
both sides forces $\Lambda_d\equiv\Lambda_{d'}\pmod{2^{\nu+1}}$,
contradicting $\vv_2(\Lambda_d-\Lambda_{d'})=\nu$. Hence
$\ord_{d'}s\le\nu$. (Small-index templates have not converged to
$\Lambda_{d'}$; this affects which small values occur, not the bound.)
\end{proof}

\begin{remark}[the family method and the price list]\label{rem:price}
Historically, Theorems A--C were first obtained for subclasses by a
family method that yields quantitative structure the shadow does not.
The \emph{burn} family $3^{j}2^{t}-1$ ascends at ratio $>\tfrac32$ per
step with frozen residue $-1$ and $\lambda\equiv\bot$, forcing $g$ to
grow at slope $\log_2\tfrac32$ along a coordinate line; the tower's
$e{=}2$ chains cap the same values from above, each level costing one
$\log_2\tfrac43$. The resulting price list --- each detector level buys
a would-be potential exactly $\log_2\tfrac43$ of budget against an
unbounded liability --- explains \emph{why} suffix information fails,
while the shadow certificates prove \emph{that} it fails. The full
family proofs are retained in the repository.
\end{remark}

\section{Certificates from raw orbits}\label{sec:certificates}

The constraint system ``$\Phi$ nonincreasing'', restricted to any
coordinate window, is a system of difference constraints
$g(v)-g(u)\le-\log_2(f(x)/x)$ over observed steps with
$u=\mathrm{coords}(x)$, $v=\mathrm{coords}(f(x))$; it is infeasible iff
the coordinate graph contains a directed cycle whose product of
value ratios exceeds $1$, and any such cycle is verifiable in exact
rational arithmetic. Mining by Bellman--Ford over odd $x\le10^{6}$
yields short certificates in every window we tested: the pair
$57\mapsto43$, $27\mapsto41$ closes a $2$-cycle in
$(x\bmod16,\tau)$-space with product $\tfrac{1763}{1539}>1$, refuting
Theorem~\ref{thm:A}'s class at $m\le4$ in two lines; a $5$-cycle
through $41\mapsto31=2^{5}-1$ does the same with the level-$1$ detector
added. Inspection of mined witnesses at larger moduli revealed the
$-5,-7$ residue pattern of Section~\ref{sec:shadow}. This pipeline is
logically independent of every lemma in this paper --- it uses raw
orbit steps only --- and would have refuted the theorems had a valid
correction existed on the tested windows.

\section{The boundary}\label{sec:boundary}

What survives Theorem~\ref{thm:B} must see $\log_2x$ to a precision
finer than every $\|a\beta\|$ --- integer and fractional parts jointly:
the \emph{quantized full logarithm}. There, a closed certificate needs
total gain in $(0,2^{-j})$, which pure shadows cannot supply
($\beta>2^{-j}$ for $j\ge3$); any certificate must mix gain and loss
segments, a connectivity question in the coordinate graph, minable per
window but with no uniform construction known.

\begin{question}
Is there a nonincreasing potential of the form
$\log_2x+g\bigl(\lfloor2^{j}\log_2x\rfloor,\ \tau(x),\
x\bmod2^m\bigr)$ for some $j,m$? As $j\to\infty$ the class approaches
arbitrary corrections, for which the existence of a monotone potential
is conjecture-adjacent; the one-step no-go program has plausibly
reached the point where further progress and the conjecture itself
interact.
\end{question}

\begin{question}
Do multi-step potentials (functions nonincreasing along $f^{k}$ for
fixed $k$, or path-dependent accountings) admit analogous obstructions?
The shadow supplies candidate constraints, since its $k$-step behaviour
is exactly known; note that any $k$-step potential must in particular
survive $k$ dividing $2a$ along the iterated shadow.
\end{question}

We also note what the results do \emph{not} say: nothing here bears on
individual trajectories, densities of divergent orbits, or cycles. The
theorems delimit a proof \emph{method}, not the dynamics.

\section{Computational verification}\label{sec:verification}

Every numbered claim is verified by scripts using exact integer and
rational arithmetic (no floating point in any assertion): the shadow
lemma over wide parameter ranges with detector evaluation at all nodes
(\texttt{verify\_shadow.py}); the iterated shadow, the concrete
Dirichlet instance $(m,j,a)=(6,4,53)$ at depth $N=177$, and the
irrationality arithmetic (\texttt{verify\_leading\_digit.py}); the
tower in all parts, including both directions of the domain criterion
and the carry-free normal form
(\texttt{verify\_tower.py}); the family-method subclasses
(\texttt{verify\_no\_local\_potential.py}, \texttt{verify\_nlp2.py});
and the independent certificate miner with exact cycle confirmation
(\texttt{verify\_nogo\_certificate.py},
\texttt{verify\_fuel\_fraction.py}). The scripts accompany the
repository.

\begin{thebibliography}{99}

\bibitem{Andaloro2000}
P.~Andaloro,
\emph{On total stopping times under $3x+1$ iteration},
Fibonacci Quart. \textbf{38} (2000), 73--78.

\bibitem{Barina2021}
D.~Barina,
\emph{Convergence verification of the Collatz problem},
J. Supercomputing \textbf{77} (2021), 2681--2688.

\bibitem{Everett1977}
C.~J. Everett,
\emph{Iteration of the number-theoretic function $f(2n)=n$,
$f(2n+1)=3n+2$},
Adv. Math. \textbf{25} (1977), 42--45.

\bibitem{KrasikovLagarias2003}
I.~Krasikov and J.~C. Lagarias,
\emph{Bounds for the $3x+1$ problem using difference inequalities},
Acta Arith. \textbf{109} (2003), 237--258.

\bibitem{Lagarias1985}
J.~C. Lagarias,
\emph{The $3x+1$ problem and its generalizations},
Amer. Math. Monthly \textbf{92} (1985), 3--23.

\bibitem{Lagarias2010}
J.~C. Lagarias (ed.),
\emph{The Ultimate Challenge: The $3x+1$ Problem},
Amer. Math. Soc., Providence, RI, 2010.

\bibitem{LagariasBibI}
J.~C. Lagarias,
\emph{The $3x+1$ problem: an annotated bibliography (1963--1999)},
arXiv:math/0309224.

\bibitem{LagariasBibII}
J.~C. Lagarias,
\emph{The $3x+1$ problem: an annotated bibliography, II (2000--2009)},
arXiv:math/0608208.

\bibitem{SimonsDeWeger2005}
J.~Simons and B.~de Weger,
\emph{Theoretical and computational bounds for $m$-cycles of the $3n+1$
problem},
Acta Arith. \textbf{117} (2005), 51--70.

\bibitem{Steiner1977}
R.~P. Steiner,
\emph{A theorem on the Syracuse problem},
Proc. 7th Manitoba Conference on Numerical Mathematics, 1977,
553--559.

\bibitem{Tao2022}
T.~Tao,
\emph{Almost all orbits of the Collatz map attain almost bounded
values},
Forum Math. Pi \textbf{10} (2022), e12.

\bibitem{Terras1976}
R.~Terras,
\emph{A stopping time problem on the positive integers},
Acta Arith. \textbf{30} (1976), 241--252.

\bibitem{Wirsching1998}
G.~J. Wirsching,
\emph{The Dynamical System Generated by the $3n+1$ Function},
Lecture Notes in Math. 1681, Springer, Berlin, 1998.

\end{thebibliography}

\end{document}
